%%% FIELD proof-structural-equation-consequences
By Definition~\texttt{def:common-response-witness}, before selector intervention the recursive equations are
\[
C=U_C,\qquad X=f_0(C,U_X),\qquad Y=B_{C,X}.
\]
The soft intervention \(\sigma_s\) is, by its declared construction, equation replacement at \(X\) only.  Hence it leaves the equations for \(C\) and \(Y\) unchanged and gives
\[
C_{\sigma_s}=U_C=C,\qquad X_{\sigma_s}=g_s(C,S),\qquad
Y_{\sigma_s}=B_{C,g_s(C,S)}.
\]
No probabilistic independence condition is used in these equalities.

Now fix \(J\subseteq\{C,X,Y\}\) and a value assignment \(j\) on \(J\).  Hard-intervention surgery replaces precisely the equations indexed by \(J\).  Recursion in the order \(C,X,Y\) therefore yields
\[
C_j=\begin{cases}j_C,&C\in J,\\ U_C,&C\notin J,\end{cases}
\qquad
X_j=\begin{cases}j_X,&X\in J,\\ g_s(C_j,S),&X\notin J,\end{cases}
\]
and
\[
Y_j=\begin{cases}j_Y,&Y\in J,\\ B_{C_j,X_j},&Y\notin J.\end{cases}
\]
The last display proves both branches in the claimed surgical response rule.  In particular, on every branch with \(Y\notin J\), the unchanged response equation is evaluated at the surgically determined parents; on every branch with \(Y\in J\), that response equation is absent and cannot be used.  Definition~\texttt{def:macro-scm} has the same response equation after reconstruction, namely \(Y=B_{C,\delta_X(H,Z_s)}\), so the identical conclusion holds in \(\mathcal M_{H,s}\).  Thus the asserted modularity is a consequence of the displayed structural equations and surgery, rather than an additional assumption.

%%% FIELD proof-projected-construction
Because \(\mathcal C\times\mathcal X\) is finite and \(0<\mu_{cx}<1\), a common finite response-table law exists explicitly: for \(b\in\{0,1\}^{\mathcal C\times\mathcal X}\), set
\[
P(B=b)=\prod_{c\in\mathcal C}\prod_{x\in\mathcal X}
\mu_{cx}^{b_{cx}}(1-\mu_{cx})^{1-b_{cx}}.
\]
Take \(U_C\sim p\) and \(S\sim\operatorname{Unif}[0,1]\), give \(U_X\) any fixed law, and take \(U_C,U_X,B,S\) mutually independent.  Then \(E[B_{cx}\mid U_C=c]=\mu_{cx}\), and, since \(C=U_C\), Assumption~\texttt{ass:selector-exogeneity} holds.  This supplies the common law \(P_R\) in Definition~\texttt{def:common-response-witness}; it is chosen once and does not depend on the selector.

Fix a feasible selector kernel \(s\) and an exogenous realization \(u=(u_C,u_X,b,s_0)\).  The map \(\widetilde\tau\) in Definition~\texttt{def:split-projection-maps} is bijective, with inverse
\[
(c,(h,z),y)\longmapsto(c,z,y).
\]
Relabeling the equations of \(\mathcal M_s\) by this bijection gives
\[
C=U_C,\qquad H=h,\qquad Z_s=g_s(C,S),
\qquad Y=B_{C,Z_s}=B_{C,\delta_X(H,Z_s)},
\]
where the last equality uses \(\delta_X(h,z)=z\).  These are exactly the equations of \(\mathcal M_{H,s}\) in Definition~\texttt{def:macro-scm}.  This proves the constructive relabeling assertion.

We next verify the all-intervention condition directly, rather than invoking it as a premise.  Fix \(J\subseteq\{C,X,Y\}\) and an assignment \(j\in\mathcal D_J\).  By Lemma~\texttt{lem:structural-equation-consequences}, the low-level surgical solution is
\[
C_j=\begin{cases}j_C,&C\in J,\\u_C,&C\notin J,\end{cases}
\quad
X_j=\begin{cases}j_X,&X\in J,\\g_s(C_j,s_0),&X\notin J,\end{cases}
\quad
Y_j=\begin{cases}j_Y,&Y\in J,\\b_{C_j,X_j},&Y\notin J.\end{cases}
\tag{1}
\]
Under the intervention split \((j^o,j^u)=\omega(j)\) from Definition~\texttt{def:split-projection-maps}, set
\[
H_j=h,qquad
Z_{s,j}=\begin{cases}j_X,&X\in J,\\g_s(C_j,s_0),&X\notin J.\end{cases}
\tag{2}
\]
Equations (1)--(2) give \(Z_{s,j}=X_j\).  If \(Y\in J\), the high-level surgery sets its outcome to \(j_Y=Y_j\).  If \(Y\notin J\), Definition~\texttt{def:macro-scm} and (2) give
\[
B_{C_j,\delta_X(H_j,Z_{s,j})}
=B_{C_j,Z_{s,j}}
=B_{C_j,X_j}
=Y_j.
\tag{3}
\]
This covers all eight possible subsets \(J\), including every simultaneous intervention involving \(Y\): in those cases the first branch of (1), rather than (3), applies.

Let \(w_j^o=(C_j,h,Y_j)\) and \(w_j^u=(\ast,Z_{s,j},\ast)\).  By (1)--(3) and the definition of \(\delta\),
\[
\delta(w_j^o,w_j^u)=(C_j,X_j,Y_j)=W_j^{\mathcal M_s}(u).
\]
On the intervened coordinates, the definition of \(\omega\) likewise gives \(\delta(j^o,j^u)=j\).  When \(X\in J\), the residual surgery fixes \(Z_s=j_X\); when \(X\notin J\), compatibility of a nonintervened residual argument means that it satisfies its retained structural equation \(Z_s=g_s(C_j,s_0)\).  Consequently (1)--(3) also prove, for every such compatible residual argument,
\[
w_j^o=\mathcal M_{H,s,[j^o]}(u,j^u,z^u),
\]
which is precisely Definition~\texttt{def:partial-projection-condition}.  Thus \(\mathcal M_{H,s}\) is the claimed partial projection, selector by selector.

It remains to calculate the causal query.  Surgery \(\operatorname{do}(H=h)\) replaces the already constant equation for \(H\) and leaves the equations for \(C,Z_s,Y\) intact.  Therefore
\[
Y=B_{C,g_s(C,S)}
\quad\text{under }\operatorname{do}(H=h).
\]
Assumption~\texttt{ass:selector-exogeneity} gives \(S\perp(C,B)\).  The response-table condition gives \(E[B_{cx}\mid C=c]=\mu_{cx}\), and \(P\{g_s(c,S)=x\mid C=c\}=s_x(c)\).  Conditioning first on \(C\) and then on the selector event hence yields
\[
\begin{aligned}
E_{\mathcal M_{H,s}}[Y\mid\operatorname{do}(H=h)]
&=\sum_{c\in\mathcal C}p_c
  \sum_{x\in\mathcal X}
  E\!\left[B_{cx}\mathbf 1\{g_s(c,S)=x\}\mid C=c\right]\\
&=\sum_{c\in\mathcal C}p_c
  \sum_{x\in\mathcal X}\mu_{cx}s_x(c).
\end{aligned}
\tag{4}
\]
The common-response-witness definition fixes \(P_R\), \(\mathcal M_R\), and the response table before \(s\) is supplied.  Equations (1)--(4) show that changing \(s\) changes only the selector equation and the induced law of \(Z_s\), as claimed.

%%% FIELD proof-sharp-envelope
First, \(\Gamma(p,r)\) is nonempty: the product array \(\gamma^0_{cx}=p_cr_x\) is nonnegative and has row margins \(p\) and column margins \(r\).  By Definition~\texttt{def:transport-polytope}, \(\Gamma(p,r)\) is an intersection of finitely many closed half-spaces and affine hyperplanes.  Each coordinate satisfies \(0\leq\gamma_{cx}\leq1\), so the set is compact; it is convex because all of its constraints are linear.

Consider any member of \(\mathfrak A_{\mathrm{free}}\).  Put \(\gamma_{cx}=p_cs_x(c)\).  Since \(s(c)\) is a probability vector,
\[
\sum_x\gamma_{cx}=p_c\sum_xs_x(c)=p_c.
\]
Assumption~\texttt{ass:calibration} gives \(\sum_c\gamma_{cx}=r_x\).  Thus \(\gamma\in\Gamma(p,r)\), and equation (4) of Lemma~\texttt{lem:projected-construction} gives
\[
\theta_h=\sum_{c,x}\gamma_{cx}\mu_{cx}.
\tag{5}
\]
This proves one inclusion in the asserted identified set.

Conversely, fix \(\gamma\in\Gamma(p,r)\).  Positivity of every \(p_c\) permits the division
\[
s_x(c)=\frac{\gamma_{cx}}{p_c}.
\tag{6}
\]
The row constraints on \(\gamma\) show that \(s(c)\in\Delta(\mathcal X)\), and the column constraints show that (6) satisfies Assumption~\texttt{ass:calibration}.  The selector construction certified by Lemma~\texttt{lem:projected-construction} realizes this kernel while retaining its common law, response table, and low-level response SCM.  Definition~\texttt{def:free-class} therefore places the resulting tuple in \(\mathfrak A_{\mathrm{free}}\), and the lemma supplies its all-intervention projection certificate and, by (5), realizes the value \(\langle\gamma,\mu\rangle\).  This proves the reverse inclusion.

The linear map \(T(\gamma)=\langle\gamma,\mu\rangle\) sends the compact set \(\Gamma(p,r)\) to a compact set and the convex set \(\Gamma(p,r)\) to a convex subset of \(\mathbb R\).  Hence its image is a closed interval.  Its endpoints are exactly the minimum and maximum in Definition~\texttt{def:endpoints}.  More explicitly, if \(\gamma_L\) and \(\gamma_U\) attain those extrema and \(L<U\), then for any \(v\in[L,U]\),
\[
\lambda=\frac{v-L}{U-L},\qquad
\gamma_v=(1-\lambda)\gamma_L+\lambda\gamma_U
\]
satisfy \(\gamma_v\in\Gamma(p,r)\) and \(T(\gamma_v)=v\).  If \(L=U\), any feasible coupling realizes the sole value.  Thus every point is attained by the certified SCM construction and
\[
\Theta_h=\{\langle\gamma,\mu\rangle:\gamma\in\Gamma(p,r)\}=[L,U].
\]
Moreover, \(\sum_{c,x}\gamma_{cx}=1\).  Since the finite array \(\mu\) has every entry strictly between zero and one,
\[
0<\min_{c,x}\mu_{cx}\leq L\leq U\leq\max_{c,x}\mu_{cx}<1,
\]
so the envelope is strictly inside the logical Bernoulli bounds.

As an exact smallest transport stress check, take two contexts and two actions, \(p=r=(1/2,1/2)\), and
\[
\mu=\begin{pmatrix}1/4&3/4\\3/4&1/4\end{pmatrix}.
\]
Every feasible coupling is
\[
\gamma(t)=\begin{pmatrix}t&1/2-t\\1/2-t&t\end{pmatrix},
\qquad 0\leq t\leq1/2,
\]
and direct substitution gives \(T(\gamma(t))=3/4-t\), hence the attained nontrivial interval \([1/4,3/4]\).  The endpoint kernels obtained from (6), together with the common Bernoulli response table and independent seed, satisfy the structural identities already verified above.

%%% FIELD proof-flow-certificates
Construct a finite directed network with a source node \(o\), one node for each \(c\in\mathcal C\), one node for each \(x\in\mathcal X\), and a sink node \(d\).  Send one unit of flow from \(o\) to \(d\); prescribe flow \(p_c\) on \(o\to c\), allow nonnegative flow \(f_{cx}\) on \(c\to x\), and prescribe flow \(r_x\) on \(x\to d\).  Equivalently, the only optimization variables are \((f_{cx})\), and flow conservation is exactly
\[
f_{cx}\geq0,qquad
\sum_{x\in\mathcal X}f_{cx}=p_c\quad(c\in\mathcal C),
\qquad
\sum_{c\in\mathcal C}f_{cx}=r_x\quad(x\in\mathcal X).
\tag{7}
\]
Give arc \(c\to x\) cost \(\mu_{cx}\) and all other arcs cost zero.  By (7), the feasible-flow set is exactly \(\Gamma(p,r)\), so this min-cost-flow problem has value
\[
\min_{f\text{ satisfying }(7)}\sum_{c,x}f_{cx}\mu_{cx}=L.
\]
Changing the costs on \(c\to x\) to \(-\mu_{cx}\) gives a second min-cost-flow problem with value \(-U\).  The feasible set is nonempty and compact by Proposition~\texttt{prop:sharp-envelope}, so both optima are attained.  Denote returned optimal flows by \(\gamma_L\in\Gamma_L\) and \(\gamma_U\in\Gamma_U\).

Because \(p_c>0\) for every context, Definition~\texttt{def:endpoint-certificates} gives well-defined kernels
\[
s_{L,x}(c)=\frac{\gamma_{L,cx}}{p_c},
\qquad
s_{U,x}(c)=\frac{\gamma_{U,cx}}{p_c}.
\tag{8}
\]
The row constraints in (7) make each row in (8) sum to one, and the column constraints give
\[
\sum_cp_cs_{L,x}(c)=r_x
=\sum_cp_cs_{U,x}(c),
\]
so both selectors obey Assumption~\texttt{ass:calibration}.  Applying Definition~\texttt{def:endpoint-certificates} to the two flows uses the same finite response SCM and common exogenous law, but partitions the independent uniform selector seed according to the two kernels in (8).  Lemma~\texttt{lem:projected-construction} then supplies, for each returned flow, the high-level SCM and its all-intervention partial-projection identities.  Proposition~\texttt{prop:sharp-envelope} and optimality give their causal means as \(L\) and \(U\), respectively.  Thus the optimal flow matrices, their normalized rows, and the instantiated structural equations are explicit endpoint certificates.

%%% FIELD proof-point-id-iff
Write \(T(\gamma)=\langle\gamma,\mu\rangle\).  Since \(\Gamma(p,r)\) is nonempty, \(L=U\) is equivalent to \(T\) being constant on \(\Gamma(p,r)\): every feasible value lies between its minimum and maximum.

Suppose first that \(L=U\).  For each \(e\in E_+\), Definition~\texttt{def:effective-support} supplies \(\gamma^{e}\in\Gamma(p,r)\) with \(\gamma^{e}_e>0\).  Since \(E_+\) is finite and nonempty, the average
\[
\bar\gamma=\frac{1}{|E_+|}\sum_{e\in E_+}\gamma^e
\]
is feasible and is strictly positive on every edge in \(E_+\).  Consider an alternating cycle
\[
c_1,x_1,c_2,x_2,\ldots,c_q,x_q,c_1
\]
and set \(c_{q+1}=c_1\).  Define an array \(d\) by adding one on \((c_j,x_j)\), subtracting one on \((c_{j+1},x_j)\), and setting all other entries to zero.  At every context and action vertex the signed incident contributions cancel, so every row and column sum of \(d\) is zero.  Because \(\bar\gamma\) is positive on the finitely many cycle edges, there is \(\varepsilon>0\) such that both \(\bar\gamma+\varepsilon d\) and \(\bar\gamma-\varepsilon d\) are nonnegative.  They retain the margins and hence belong to \(\Gamma(p,r)\).  Constancy of \(T\) now implies
\[
0=\langle d,\mu\rangle
=\sum_{j=1}^q\bigl(\mu_{c_jx_j}-\mu_{c_{j+1}x_j}\bigr),
\]
which is the alternating-cycle condition.

Assume next that every such cycle contrast is zero.  Work separately on each connected component of the bipartite graph with vertex set \(\mathcal C\sqcup\mathcal X\) and edge set \(E_+\).  Choose a spanning tree in the component, fix one context-root potential to be zero, and recursively define potentials along tree edges by
\[
b_x=\mu_{cx}-a_c\quad\text{when }a_c\text{ is already defined},
\qquad
a_c=\mu_{cx}-b_x\quad\text{when }b_x\text{ is already defined}.
\tag{9}
\]
The unique tree path makes (9) well defined and ensures \(\mu_{cx}=a_c+b_x\) on every tree edge.  If \((c,x)\) is a non-tree edge, adjoining it to the tree creates one alternating cycle.  Substitution of the tree equalities into that cycle's zero contrast telescopes to
\[
\mu_{cx}-a_c-b_x=0.
\]
Thus additive potentials exist on every effective edge of every component.

Finally, suppose these potentials exist.  If \(\gamma_{cx}>0\) for a feasible \(\gamma\), then \((c,x)\in E_+\) by the definition of effective support.  Consequently
\[
\begin{aligned}
T(\gamma)
&=\sum_{(c,x)\in E_+}\gamma_{cx}(a_c+b_x)\\
&=\sum_c a_c\sum_x\gamma_{cx}+\sum_xb_x\sum_c\gamma_{cx}\\
&=\sum_cp_ca_c+\sum_xr_xb_x,
\end{aligned}
\tag{10}
\]
which does not depend on \(\gamma\).  Hence \(L=U\).  This proves all three equivalences.  Under the present strictly positive margins, the product coupling \(p_cr_x\) is positive on every pair, so \(E_+=\mathcal C\times\mathcal X\); the support-aware proof above also shows exactly how the criterion restricts if structural zero edges are imposed.

%%% FIELD proof-fixed-map-obstruction
Fix \(q_{\mathrm{can}}\).  Under context law \(p\), its joint context--residual law is necessarily
\[
P(C=c,Z_s=x\mid\operatorname{do}(H=h))
=p_cq_{\mathrm{can}}(x\mid c,h)
=\gamma_{\mathrm{can},cx}.
\tag{11}
\]
Thus a fixed canonical residual kernel realizes the singleton set \(\{\gamma_{\mathrm{can}}\}\).  It spans the feasible coupling set precisely when
\[
\Gamma(p,r)=\{\gamma_{\mathrm{can}}\};
\]
in particular, this equality itself entails that (11) has column margins \(r\).

Let \(m=|\mathcal C|\) and \(k=|\mathcal X|\).  The affine equations defining \(\Gamma(p,r)\) have one redundancy because both the row totals and the column totals sum to one.  More explicitly, in their real affine solution set, the \((m-1)(k-1)\) entries \(\gamma_{cx}\) with \(c<m\) and \(x<k\) may be chosen freely; the remaining entries are then uniquely recovered from
\[
\gamma_{c k}=p_c-\sum_{x<k}\gamma_{cx},
\qquad
\gamma_{m x}=r_x-\sum_{c<m}\gamma_{cx},
\]
with \(\gamma_{mk}\) fixed consistently by either final margin equation.  Hence that affine solution set has dimension \((m-1)(k-1)\).  The product coupling \(\gamma^0_{cx}=p_cr_x\) is strictly positive because all margins are positive.  A sufficiently small relative neighborhood of \(\gamma^0\) in the affine solution set remains nonnegative, so the transportation polytope has the same dimension.  If \(m,k\geq2\), this dimension is positive; therefore \(\Gamma(p,r)\) is not a singleton and cannot be spanned by (11).

For a direct two-context, three-action check, take
\[
p=(1/2,1/2),\qquad r=(1/3,1/3,1/3),\qquad
\mu=\begin{pmatrix}1/10&1/2&9/10\\9/10&1/2&1/10\end{pmatrix}.
\]
The two feasible endpoint couplings are
\[
\gamma_L=\begin{pmatrix}1/3&1/6&0\\0&1/6&1/3\end{pmatrix},
\qquad
\gamma_U=\begin{pmatrix}0&1/6&1/3\\1/3&1/6&0\end{pmatrix}.
\]
Indeed, writing the first-row flows as \(a_x\) reduces the objective, up to a constant, to \(-\frac45a_1+\frac45a_3\) subject to \(0\leq a_x\leq1/3\) and \(\sum_xa_x=1/2\), which proves the displayed minimizer and maximizer.  Their causal means are \(7/30\) and \(23/30\).  Normalization by \(p_c\) gives two distinct selector kernels, and Lemma~\texttt{lem:projected-construction} verifies every partial-projection surgery for both while retaining the common response SCM.  By contrast, the fixed independent canonical kernel \(q_{\mathrm{can}}(x\mid c,h)=r_x\) gives the single product coupling with every entry \(1/6\) and causal mean \(1/2\).  It cannot equal both endpoint couplings.  This finite check confirms the quantifier boundary: Definition~\texttt{def:free-class} permits a selector-specific residual law, and the construction proves Definition~6 compatibility separately for each selector; it does not preserve one fixed canonical Equation-11 Layer-3 map across the family.

%%% FIELD proof-additive-reduction
Assume that on every effective edge there are potentials satisfying \(\mu_{cx}=a_c+b_x\).  For any selector in \(\mathfrak A_{\mathrm{free}}\), Proposition~\texttt{prop:sharp-envelope} associates the coupling \(\gamma_{cx}=p_cs_x(c)\in\Gamma(p,r)\) and gives \(\theta_h=\sum_{c,x}\gamma_{cx}\mu_{cx}\).  A positive entry of \(\gamma\) is, by Definition~\texttt{def:effective-support}, an edge of \(E_+\).  Terms off \(E_+\) therefore have zero weight, and the margin equations yield
\[
\begin{aligned}
\theta_h
&=\sum_{(c,x)\in E_+}\gamma_{cx}(a_c+b_x)\\
&=\sum_ca_c\sum_x\gamma_{cx}+\sum_xb_x\sum_c\gamma_{cx}\\
&=\sum_cp_ca_c+\sum_xr_xb_x.
\end{aligned}
\]
The right-hand side contains no selector kernel.  This also follows from the potential characterization in Lemma~\texttt{lem:point-id-iff}, which ensures \(L=U\).

Under full effective support, the additive representation holds on all of \(\mathcal C\times\mathcal X\).  Any two calibrated conditional version laws \(s\) and \(s'\) have the same margins \(p,r\), so applying the displayed calculation to each gives
\[
\theta_h(s)=\sum_cp_ca_c+\sum_xr_xb_x=\theta_h(s').
\]
Thus changing the calibrated conditional version law cannot change the macro mean, which is exactly the asserted version-irrelevance reduction.

%%% FIELD proof-uniform-superiority
By Definition~\texttt{def:control-contrast} and Proposition~\texttt{prop:sharp-envelope}, the set of effects relative to the reference mean is
\[
\Theta_h-\theta_0=[L-\theta_0,U-\theta_0].
\]
The lower endpoint is attained by a compatible selector SCM.  Therefore every compatible selector is strictly superior precisely when every element of this interval is positive, which holds if and only if its attained minimum satisfies
\[
L-\theta_0>0,
\]
or equivalently \(L>\theta_0\).  Likewise, every compatible selector is weakly noninferior precisely when every element is nonnegative, which is equivalent to \(L-\theta_0\geq0\), or \(L\geq\theta_0\).  Endpoint attainment rules out replacing either strict or weak inequality by a weaker condition.
